\begin{definition}
    $\tqbf = \{ \varphi\ |\ \exists x_1 \forall x_2 \exists x_3 \ldots (\exists / \forall) x_n \ \varphi(x_1,...,x_n) = 1 \}$
\end{definition}

\begin{theorem}
    $\tqbf$ — \psp-полный.
\end{theorem}
\begin{proof}
\begin{enumerate}
    \item $\tqbf \in \psp$.
    
    Если значения всех переменных заданы, то формула вычисляется на полиномиальной доппамяти (при должном упорстве на константной, но это не необязательно). Алгоритм будет заключаться в рекурсивном переборе всех возможных значений. Допустим мы умеем вычислять истинность $\exists x_{i + 1}, \forall x_{i + 2}\ldots \varphi(x_1, \ldots, x_i, x_{i + 1}, \ldots)$ для некоторого четного $i$ и произвольных $x_1, \ldots, x_i$ или $\forall x_{i + 1}, \exists x_{i + 2}\ldots \varphi(x_1, \ldots, x_i, x_{i + 1}, \ldots)$ для некоторого нечетного $i$ и произвольных $x_1, \ldots, x_i$. Тогда, совершив два рекурсивных вызова и взяв их конъюнкцию (дизъюнкцию) результатов научимся вычислять истинность $\forall x_i \exists x_{i + 1}, \forall x_{i + 2}\ldots \varphi(x_1, \ldots, x_i, x_{i + 1}, \ldots)$ ($\exists x_i \forall x_{i + 1}, \exists x_{i + 2}\ldots \varphi(x_1, \ldots, x_i, x_{i + 1}, \ldots)$) для произвольных $x_1, \ldots, x_{i - 1}$. Тогда на каждом уровне рекурсии мы тратим полином памяти, а всего таких уровней $n$.
    
    \begin{lemmanote}
        Достаточно константной памяти на каждом уровне и линейной на все-все-все.
    \end{lemmanote}
    
    \item $\tqbf$ — $\mathbf{PSPACE}$-сложен. 
    
    \begin{lemmanote}
        Доказательство из compl-book имеет право на существование, но это какая-то уж полная дичь % осуждаю
        (там нужно пруфать теорему Сэвича% охуеть, надо что-то доказывать
        , вводить $\mathsf{SUCCINPATH}$, а он страшен и доказывать сначала его полноту%трудности достаточно, держу в курсе
        ), поэтому берём из вики.
        % МУСАТОВ ПОМОГИ МУСАТОВ ПОМОГИ МУСАТОВ ПОМОГИ
        % пиздец...

    \end{lemmanote} 
    
    Пусть $L \in \psp$, то есть $L$ разрешим детерминитованной МТ на полиномиальной памяти. Заметим, что каждую такую машину можно описать булевой формулой, где переменные отвечают за ее состояния, содержимое лент и положения головок. Наше сведение будет использовать переменные $c_1$, $c_2$ и натуральное число $t$, которые представляют две возможные конфигурации МТ для $L$. Будем строить формулу $\varphi(c_1, c_2, t)$, которая будет выполнима тогда и только тогда, когда от $c_1$ до $c_2$ можно дойти за не более чем $t$ шагов. Тогда наша сводящая функция (вспомним, что она полиномиальна по времени и, как следствие, по памяти, так как мы хотим именно этого в определение \psp сводимости) построит $\varphi(c_{start}, c_{accept}, T)$ и $T$~--- максимальное число шагов, которое требуется, чтобы дойти от одного состояния до другого. $T = \mathcal{O}(\exp(n))$, так как в $\psp$ память ограничена полиномом и, как следствие, число всех возможных ее состояний экспоненциально. Более того, мы еще не разобрались с потенциальными циклами, но тогда мы никогда не достигнем принимающего состояния никогда (а значит слово не из $L$ и проблем нет).
    
    Заметим, что нам осталось узнать, а можно ли построить эту булеву формулу за полиномиальное время, то есть является ли сведение сведением по Карпу? Для ответа на этот вопрос построим явно формулу по индукции.
    
    \begin{itemize}
        \item $t = 1$. Так как МТ детерминированная, просто надо проверить, что на всех лентах содержимое не изменилось кроме как в положениях, куда указывали головки. Построимо за линейное время от длины входа (или всех лент).
        \item $t > 1$. Тут мы будем делать аналог meet-in-the-middle. Запишем нашу формулу рекурсивно:
        $$
        \varphi(c_1, c_2, t) = \exists m\brackets{\varphi(c_1, m, t/2) \wedge \varphi(m, c_2, t/2)}
        $$
    \end{itemize}
    
    Теперь заметим, что $t$ ограничено только $T$, а оно экспоненциально по длине входа. Более того, можно получить таким образом, что время построения экспоненциально (так как длина будет экспоненциальной путем решения рекурренты $l(n) = 2\cdot l(n/2)$), а это неприемлимо. Давайте сокращать формулу, для этого заведем переменные $c_3$ и $c_4$, отвечающие конфигурациям $(c_1, m)$ и $(m, c_2)$. Тогда формулу выше можно переписать:
    $$
    \varphi(c_1, c_2, t) = \exists m\forall (c_3, c_4) \in \{(c_1, m), (m, c_2)\} \brackets{\varphi(c_3, c_4, t/2)}
    $$
    Получаем, что длина формулы равна длине рекурсивной цепочки (точнее $l(n) = \mathcal{O}(1) + l(n/2)$), а она уже полиномиальна (логарифм от экспоненты). А раз такая формула построима за полиномиальное время, то у нас есть сведение, так как принадлежность слова языку $L$ равносильна тому, что построенная формула выполнима.
\end{enumerate}

\end{proof}